\documentclass[11pt]{article}

\usepackage[margin=1in]{geometry}
\usepackage{amsmath,amssymb}
\usepackage{booktabs}
\usepackage{multirow}
\usepackage{hyperref}
\usepackage{microtype}
\usepackage{enumitem}

\hypersetup{
  colorlinks=true,
  linkcolor=black,
  citecolor=black,
  urlcolor=blue
}

\title{Adaptive Temporal Retrieval\\
for Time-Sensitive Question Answering}

\author{
Shashank Kumar (08619051623) \quad Daksh Agarwal (10819051623)\\
University School of Automation and Robotics\\
Guru Gobind Singh Indraprastha University\\[0.35em]
{\normalsize Supervisor: Dr.\ Annu Priya}\\[0.15em]
{\normalsize ARP-455 minor project}
}

\date{November 2026}

\begin{document}
\maketitle

\begin{abstract}
The student project behind this paper is called ``When and What to Retrieve.''
The pipeline can label a question as present, past, or neither, keep a date window around the asked time, and read a sentence whose date agrees with the question.
This paper reports what that pipeline actually changes, and where the measurements stop.

On a templated test of 77 questions, with a shared reader and shared projection, Adaptive and a fixed $\pm 2$ year filter both contain the gold string on 76 questions ($0.987$).
Exact match is $0.948$ against $0.935$: one extra question.
Precision at 5 is higher for Adaptive ($0.668$ against $0.608$).
Skipping retrieval fires on 10 of those 77 questions.
Mean latency does not fall ($1.690$\,s against $1.555$\,s).
On 193 questions about six entity names, Adaptive is ahead of the same fixed window on exact match ($0.762$ against $0.684$) and containment ($0.762$ against $0.710$), and it always retrieves.
The extra work shows up as mean latency ($2.584$\,s against $1.546$\,s).
On the full 500-question TimeQA test, Adaptive containment is $0.296$, against $0.270$ for the fixed window and $0.260$ for standard retrieval.
Exact match stays near zero for every system ($0.008$, $0.010$, $0.006$) even after span extraction.
Recall at 5 is higher for Adaptive ($0.488$) than for B3 ($0.392$), which is the cost of filtering on stored timestamps that are often not fact dates.
B4, a TempRALM-style dense+BM25 ranker with exponential year decay and no hard cutoff, scores exact match $0.662$ and containment $0.727$ on that templated test, well behind Adaptive.
The sufficiency loop does not raise answers: turning it off leaves exact match unchanged and, on TimeQA, containment moves from $0.296$ to $0.298$.
\end{abstract}

\section{Introduction}
\label{sec:intro}

A fact such as who holds an office is true on an interval.
Retrieving the nearest passage can return a true sentence about another year.
Retrieval-augmented generation grounds a generator in an external corpus \cite{lewis2020rag}, but the usual pipeline has no separate decision for ``now'' and ``in 2016''.

This project adds a temporal policy on top of one dense index:
a classifier with a two-way necessity head and a three-way time head;
a date window whose width depends on the cue;
a sufficiency loop that may widen the search;
and a reader that prefers a sentence whose date matches the question.
The project title asks two questions, when to retrieve and what to retrieve.
The second is what the code mostly implements.
The first is a whitelist.
Section~\ref{sec:when} says exactly when retrieval is skipped, and Section~\ref{sec:results} shows that this skip does not change latency and does not occur on the held-out or TimeQA runs.

We do not claim a new state of the art.
We do not claim that the adaptive window beats a fixed window once the answer stage is equal, except where a table shows that directly.
One frozen Checkpoint~B configuration is reported on the full synthetic, held-out, and TimeQA tests.

The contribution is a description of that pipeline, the formulas it uses, and a reading of the saved numbers that does not pick the friendliest metric.
The frozen comparison in Section~\ref{sec:results} uses one Checkpoint~B configuration on the full synthetic, held-out, and TimeQA tests, with a shared reader and shared answer projection.
What remains missing --- a published temporal retriever as a fifth row, a live clock, and a learned skip that is not a trivia whitelist --- is listed in Section~\ref{sec:limits}.

\section{Related work}
\label{sec:related}

\paragraph{Retrieval-augmented generation.}
Lewis et al.\ \cite{lewis2020rag} condition a generator on passages from a dense index.
Dense passage retrieval showed that a dual encoder can replace a sparse first stage for open-domain questions \cite{karpukhin2020dpr}.
Atlas trains the retriever and the reader together \cite{izacard2023atlas}.
Self-RAG learns when a generation should retrieve and when it should critique the result \cite{asai2024selfrag}.
Adaptive-RAG chooses among no retrieval, one retrieval, and several retrievals according to question complexity \cite{jeong2024adaptiverag}.
That is the closest system in name.
It is not a temporal ranker, and we do not run it.
A reader should not treat the shared word ``adaptive'' as a comparison.

\paragraph{Time in retrieval, before RAG.}
Li and Croft put document age into a language model so that recency-seeking queries do not depend on a hand-tuned boost \cite{li2003time}.
Kanhabua and N{\o}rv{\aa}g compared mixture models that add a temporal similarity to text similarity, on a news archive with time-scoped judgments \cite{kanhabua2011ranking}.
Those papers already treat ``what time is relevant?'' as a ranking problem.
Our window is a heuristic descendant of that idea, not a replacement of those models, and we do not evaluate on their collections.

\paragraph{Time in question answering.}
TempQuestions evaluates time expressions against a knowledge graph \cite{jia2018tempquestions}.
TimeQA aligns time-evolving Wikidata facts with Wikipedia \cite{chen2021timeqa}.
The metric in that paper is exact match.
Fusion-in-Decoder fine-tuned on Natural Questions and TimeQA reaches 60.5 exact match on the easy test and 46.8 on the hard test; human exact match on the hard test is 87.0 \cite{chen2021timeqa}.
The abstract's ``about 46\%'' is the hard-test figure, not a single number for the whole benchmark.
SituatedQA shows that one question can have different answers at different times \cite{zhang2021situatedqa}.
StreamingQA and RealTime QA test adaptation as new text arrives \cite{liska2022streamingqa,kasai2023realtimeqa}.
Dhingra et al.\ condition a language model on a timestamp so that closed-book recall of a changing fact depends on the year \cite{dhingra2022talm}.
FreshQA asks models about facts that change quickly, including false premises \cite{vu2023freshqa}.
TempRALM re-ranks retrieved passages with a temporal score and masks passages dated after the query time \cite{gade2024tempralm}.
We cite these systems to locate the project.
We did not reimplement them, and our TimeQA subset is not comparable to their reported scores.

\paragraph{Readers and the first-stage stack.}
SQuAD 2.0 requires a span and an abstention decision \cite{rajpurkar2018squad2}.
The learned path uses a frozen SQuAD 2.0 model and does not fine-tune it.
Embeddings are BGE-small \cite{xiao2024cpack}, searched with an exact FAISS inner-product index \cite{johnson2019faiss}.
A MiniLM cross-encoder trained on MS MARCO reorders the first-stage hits \cite{wang2020minilm,nguyen2016msmarco,reimers2019sbert}.
The router is DistilBERT \cite{sanh2019distilbert,devlin2019bert}.

\section{Problem setup}
\label{sec:setup}

The corpus is $\mathcal{D}=\{(d_i,\tau_i)\}$.
Here $\tau_i$ is the timestamp stored in the index.
It is not the only date in $d_i$.
On the templated corpus, $\tau_i$ is the date of the fact.
On the TimeQA subset, $\tau_i$ is the first date the tagger finds in the chunk, or 1 January 1970 when it finds none.
A question $q$ may also name a target interval.
The system returns a string $\hat{y}$ and a bit saying whether retrieval ran.

Answer strings are normalised as in SQuAD: lower case, punctuation removed, whitespace collapsed.
Exact match is $1$ when the strings are equal.
Token F1 is the harmonic mean of token precision and token recall, taking the best gold string when several are listed.
Containment is
\begin{equation}
\mathrm{Cont}(\hat{y}, y) =
\begin{cases}
1 & \text{if } y \subseteq \hat{y} \text{ or } \hat{y} \subseteq y,\\
0 & \text{otherwise.}
\end{cases}
\end{equation}
The second branch matters.
A short prediction inside a long gold string scores $1$, and so does a 400-character sentence that merely contains the gold span.
On TimeQA the heuristic answers are sentences, so containment mostly asks whether the selected sentence contains the gold text.
It is not an answer-quality metric.
We report exact match beside it for that reason.

If the question is marked unanswerable, or the gold string is empty, exact match, token F1, and containment are replaced by abstention accuracy:
$1$ when the system abstains on an unanswerable question, and $1$ when it does not abstain on an answerable one.
Abstention means an empty string or a short refusal such as ``insufficient evidence''.
A score of $0.250$ on 16 questions means the decision was right on 4 of them.

Precision and recall at $k{=}5$ use the gold document ids stored with the question.
Precision divides the number of gold ids in the retrieved list by the length of that list, capped at 5.
On the templated data the generator stores up to three passages of the named entity, preferring those nearest the asked date.
On TimeQA a gold id is a chunk that contains an answer string, and the list is capped at 8.
A precision of $0.673$ is therefore an average hit rate inside the top 5, not evidence that each question has three or four relevant documents.
Questions with an empty gold list are omitted from the average.

Temporal accuracy looks only at $\tau$ of the top document.
For a current question, that timestamp must fall within three years of 1 October 2025.
For a historical question, its year must fall within two of a year written in the question or the time hint.
The metric does not read $\hat{y}$.
It is undefined for TimeQA and is left blank there.
For the fixed-window baseline the metric restates the filter.
That baseline keeps documents whose stored year is within $\pm 2$ of the hinted year, and the metric asks whether the top stored year is within $\pm 2$.
If the filter keeps any document, every survivor passes, and so does the top document.
The score can fall below $1$ only when the filter is empty and the fallback document lies outside the band.
It is not independent evidence that the answer is about the right year.

\section{Method}
\label{sec:method}

The pipeline classifies the question, retrieves with the dense index, reorders with the cross-encoder, applies a temporal policy, checks sufficiency, and reads an answer.
The index, the cross-encoder, the extractive reader, answer projection, and the entity or surname cut are shared across B1--B3 and Adaptive in the frozen comparison.
The systems differ on whether retrieval can be skipped, how the time window is chosen, and whether a sufficiency loop may retrieve again.
Section~\ref{sec:exp} states that split before any table.

\subsection{Query router}
\label{sec:router}

DistilBERT maps $q$ to the classification-token state $h$.
One linear head emits necessity logits $z^{\mathrm{nec}}\in\mathbb{R}^{2}$.
Another emits temporal logits $z^{\mathrm{tmp}}\in\mathbb{R}^{3}$, over current, historical, and timeless.
A third head emits a confidence logit $z^{\mathrm{conf}}$.
Training uses class weights $w$ and minimises
\begin{equation}
\mathcal{L}
= \mathrm{CE}\!\left(z^{\mathrm{nec}}, y^{\mathrm{nec}}; w^{\mathrm{nec}}\right)
+ \lambda\,
\mathrm{CE}\!\left(z^{\mathrm{tmp}}, y^{\mathrm{tmp}}; w^{\mathrm{tmp}}\right),
\label{eq:loss}
\end{equation}
with $\lambda = 1.75$.
The confidence head is not in $\mathcal{L}$.
Nothing trains its weights, so they stay at initialization.
The encoder still moves, because the two classification losses depend on it.

At inference the class logits are divided by a temperature $T$ before the softmax.
Let $p^{\mathrm{nec}}$ and $p^{\mathrm{tmp}}$ be those distributions, and let $c=\sigma(z^{\mathrm{conf}})$.
The scalar that the policy reads is
\begin{equation}
\gamma =
\begin{cases}
\tfrac12 c + \tfrac12 \max\{\max p^{\mathrm{nec}}, \max p^{\mathrm{tmp}}\}
& \text{when blending is on,}\\[0.3em]
\max\{\max p^{\mathrm{nec}}, \max p^{\mathrm{tmp}}\}
& \text{otherwise.}
\end{cases}
\label{eq:conf}
\end{equation}
The main analyzer blends.
Because $c$ was never trained, half of $\gamma$ is the output of an untrained sigmoid.
If $\gamma < 0.55$, the policy does not trust the temporal class and retrieves in semantic mode.
An untrained head can therefore change which ranker runs.
We did not count how often $\gamma$ falls below $0.55$ on the test sets.
The router used on the learned path turns blending off, so its $\gamma$ is the maximum class probability.

\subsection{When retrieval is skipped}
\label{sec:when}

Retrieval is skipped only when the necessity class is $0$ and the question string contains one of a fixed list of trivia cues: arithmetic such as ``2 + 2'', ``H2O'', ``days are in a week'', ``spell '', and a few similar items.
A necessity prediction of $0$ on any other question still retrieves, under the name of a soft timeless lookup.
Low confidence also retrieves.
The skip is a whitelist guarded by a classifier, not a learned abstention from retrieval.

There is no parametric model underneath the skip.
The closed-book control is the same extractive reader with an empty document list, plus a few hardcoded replies for those trivia items.
A skip cannot be scored as ``the model already knew the answer.''
It can only be scored as ``the index was not called.''

\subsection{First-stage retrieval}

Documents are embedded with BGE-small English v1.5.
Queries use that model's required prefix.
Search is FAISS \texttt{IndexFlatIP}, exact on these corpora.
Evaluation asks for 80 candidates.
Their FAISS scores are min--max normalised to $\tilde{s}^{\mathrm{f}}_i\in[0,1]$ inside the pool.
The cross-encoder score $s^{\mathrm{ce}}_i$ is left in its raw range.
The blend is
\begin{equation}
s_i = 0.35\,\tilde{s}^{\mathrm{f}}_i + 0.65\,s^{\mathrm{ce}}_i .
\label{eq:blend}
\end{equation}
The top 40 values of $s_i$ are kept.
The two terms do not share a scale, so $0.35$ and $0.65$ are not the weights of a mixture of comparable quantities.
A cross-encoder score of $6$ dominates a dense score in $[0,1]$ whatever the written weights say.
We did not normalise $s^{\mathrm{ce}}$ in the reported runs, and we did not choose the weights on a development split.
They are constants from implementation.

\subsection{Dates}
\label{sec:dates}

Dates are read by rules.
ISO dates and month-year expressions are kept.
A bare year is placed on 15 June, and that midpoint is dropped when the same year already has a fuller date.
Intervals cover $y$--$y$, ``from \ldots to \ldots'', ``until'' or ``through'', two-digit seasons such as 2004--05, and ``for $n$ years beginning in $y$'' with $n\le 40$.
A span longer than 45 years is discarded, which is aimed at birth--death pairs and will also discard a long reign or a long career.
``(born \ldots)'' in the question is stripped before the target year is chosen.
``Before month year'' becomes the interval from three years earlier through that month.
Relative phrases such as ``last year'', ``since 2010'', and ``the predecessor'' are not given their own target.
A decade cue is recognised only when the tagger emits one; there is no general ordinal logic.

Let $I(d)$ be the intervals in a passage and $U(d)$ the point dates, including a usable timestamp.
The distance from a target date $t$ to the passage is
\begin{equation}
\delta(d, t) =
\begin{cases}
0 & \text{if } t \in [u,v] \text{ for some } [u,v]\in I(d),\\[0.25em]
\min\limits_{x \in \partial I(d)\cup U(d)} |x-t| & \text{otherwise,}
\end{cases}
\label{eq:delta}
\end{equation}
where $\partial I(d)$ is the set of endpoints.
Distance zero means the asked date falls inside a stated range.

\subsection{Window and subject filter}

Historical questions use the half-widths in Table~\ref{tab:windows}.
While no surviving document satisfies $\delta(d,t)\le 365.25\,h$, the half-width grows by
\begin{equation}
h \leftarrow \min(h_{\max},\; 1.5h + 0.25).
\label{eq:grow}
\end{equation}
Growth stops at the first in-window survivor or at $h_{\max}$.
We did not ablate the early stop, so we cannot say that stopping early changes the answer.
The widths themselves were set by hand.

\begin{table}[t]
\centering
\caption{Initial and maximum half-width of the historical window, in years.}
\label{tab:windows}
\begin{tabular}{lcc}
\toprule
Cue & Initial $h$ & Maximum $h$ \\
\midrule
Exact calendar date & 0.35 & 1.5 \\
Early, mid, or late in a year & 1.5 & 3.5 \\
A bare year & 1.0 & 4.0 \\
A decade & 5.0 & 8.0 \\
Historical, but no year & 3.0 & 8.0 \\
\bottomrule
\end{tabular}
\end{table}

On the historical branch, before the window test, the pool is cut to a subject.
If any candidate has a non-empty entity field, only documents whose overlap with that field is within $0.01$ of the best overlap are kept.
In the frozen comparison this cut is also applied to B2 and B3, so a precision gap is not an entity-field exclusive.
On the templated corpus the field is the entity name, filled when the passage was generated.
Current mode does not use this cut.
It keeps a lexical tie and sorts by timestamp, and the fixed-window baseline calls that same function when it finds no year.
The overlap text includes the entity string whenever the field is filled, so both systems see it there.
If the field is empty, as it is for every TimeQA chunk, a document is kept when lexical overlap is at least $0.05$ and the text contains a surname, unless overlap is within $0.01$ of the best document.
The surname is the last capitalised token of the first person-like span, ignoring months and question words, and it must be at least three characters.
The rule is English-specific.
It drops a pronoun or a given name, and it was added while inspecting TimeQA failures.

After a stop list, if $H$ remaining query tokens occur in the document, the lexical boost is $b=\min(0.25, 0.05H)$.
The stop list includes months, years, and words such as ``between'', ``held'', and ``team'', because those tokens were matching the wrong passage.

\subsection{Temporal score}

Among documents that survive, semantic scores are min--max normalised to $\tilde{s}_i$.
For current questions, with reference date $r=\text{1 October 2025}$,
\begin{equation}
g^{\mathrm{cur}}_i = \max\left(0,\; 1 - \frac{\max(r-\tau_i,\, 0)}{3650}\right).
\label{eq:cur}
\end{equation}
The denominator is about ten years.
Current mode then keeps documents whose lexical overlap is within $0.01$ of the best and sorts them by $\tau_i$, latest first.
The displayed score of rank $i$ in that tie is $1-0.05\,i/n$, so an older tied document is not treated as weak evidence by the sufficiency check.
``Latest timestamp'' is a proxy for ``current holder'' only when the index date is the fact date.
The reference day is frozen, and it is more than a year before the date of this paper.
No live benchmark is used.

For historical questions,
\begin{equation}
g^{\mathrm{hist}}_i
= \max\left(0,\;
1 - \frac{\delta(d_i, t)}{\max(90,\; 365.25\,w)}
\right),
\label{eq:hist}
\end{equation}
with $w$ the window width in years.
On a tagged corpus, and only when the caller sets a time-first flag, rank is
\begin{equation}
f_i = g^{\mathrm{hist}}_i + 0.01\,\tilde{s}_i + b_i .
\label{eq:timefirst}
\end{equation}
The fixed-window baseline does not set that flag.
On untagged prose with the flag set, a shrunk semantic score is avoided by $\tilde{s}'_i=0.5+0.5\tilde{s}_i$ and
\begin{equation}
f_i = 0.45\,\tilde{s}'_i + 0.55\,g^{\mathrm{hist}}_i + b_i .
\label{eq:untagged}
\end{equation}
Otherwise
\begin{equation}
f_i = \alpha \tilde{s}_i + (1-\alpha)\,g_i\tilde{s}_i + b_i ,
\label{eq:gate}
\end{equation}
with $\alpha=0.6$.
The product $g_i\tilde{s}_i$ is the default, including for the fixed-window baseline's call into the same function.
The top $k=5$ documents go to the reader.
Tagged and untagged questions therefore do not use the same ranking formula.
Saying that one pipeline was applied to every dataset would hide that branch.

\subsection{Sufficiency}

During retrieval the judge sees the min--max semantic score of each document after temporal reordering, not the raw cross-encoder score.
Let $s_{(1)}$ be that score for the top-ranked document, $\overline{s}$ the mean, $C$ the fraction of content query tokens found in the retrieved text, $T\in[0,1]$ a temporal fit, and $b_{(1)}$ the lexical boost of the top document.
\begin{equation}
\begin{aligned}
S =\;& 0.35\,C + 0.25\,\min(1, s_{(1)}) + 0.15\,\overline{s}\\
&+ 0.15\,T + 0.10\,\min(1, b_{(1)}/0.25).
\end{aligned}
\label{eq:suff}
\end{equation}
Because the scores were min--max normalised inside the pool, $s_{(1)}=1$ whenever the temporally top document is also the semantic maximum, which is common.
The $0.25$ term is then a constant.
A pool of poor documents that still differ from each other receives the same saturated value.
$T$ takes the better of the stored timestamp and $\delta(d,t)$, so a passage whose index date is a birth year can still score well if the text contains the asked date.
Evidence is accepted during retrieval when the subject check passes and $S\ge 0.55$.
The weights and the threshold were not fit to data.

If the check fails, the loop may enlarge the candidate pool, raise $k$, add to the window, or fall back toward semantic ranking.
At most two repair rounds run, for three retrieves in total.
A wall-clock budget of 60 seconds per question stops the loop and keeps the best evidence so far.
Both stops are logged.
After an answer is drafted, an abstention stub can trigger one semantic retry only if the budget still allows it, and the retry replaces the answer only when its sufficiency score is strictly higher.

\subsection{Reading the answer}
\label{sec:reader}

The offline reader does not call an external language model.
B1, B2, B3, and Adaptive share this reader and then the same projection.

The shared reader has a few shortcuts.
Trivia cues have fixed replies.
Definition questions return the first sentence.
A passage that contains only one long sentence, and matches an office or winner template, is shortened to the name.
Initials are not sentence breaks, so ``B.\ R.\ Gavai'' stays one token.
That rule was added after a split initial hurt exact match on a templated question.

For a longer passage the reader scores sentences by date agreement, content-word overlap, and a bonus for higher-ranked documents.
A fragment shorter than 24 characters, including a section header, stays attached to the next sentence.
A dated clause shorter than 48 characters is prepended with the previous sentence.
The best sentence is truncated at 400 characters.
This selector is shared.
It was developed against TimeQA failures, including career headers and list sentences.

Projection extracts a name from general who-question and result patterns such as ``held the office'' or ``$X$ won''.
On TimeQA, evaluation also extracts a short span from the selected sentence before exact match and token F1, while containment still uses the sentence.
The patterns do not reliably extract a title such as ``assistant district attorney'' from a sentence that mentions it in passing.

\subsection{Learned path}
\label{sec:learned}

The learned path is a separate stack.
It is not the system in the full-test tables.

Its router starts from the saved analyzer, freezes the bottom four of six DistilBERT layers, and trains for one epoch on 240 questions.
Temperature $T=2.8$ was chosen on 48 questions, and the accuracy we report is on those same 48 questions: necessity accuracy $0.812$ (F1 $0.889$) and temporal accuracy $0.812$ (macro-F1 $0.782$).
That is not a held-out estimate.
The drop from the original analyzer's validation accuracy of $1.0$ is also not a like-for-like change.
The original figure is on the templated validation split.
The $0.812$ figure is on this 48-question calibration split.
The original weight file was not overwritten.

The ranker is gradient boosting \cite{friedman2001gbm}, 80 trees of depth 3, on 357 rows of which 143 are positive.
Features include the cross-encoder and FAISS scores, lexical overlap, date distance, the router's class probabilities, and two products: current-probability times recency, and historical-probability times $1/(1+\text{year distance})$.
At inference a temporal prior
\begin{equation}
\pi_i
= p_{\mathrm{cur}}\,\rho_i
+ p_{\mathrm{hist}}\frac{1}{1+\Delta_i}
+ p_{\mathrm{time}}\,\tilde{s}^{\mathrm{ce}}_i
\label{eq:prior}
\end{equation}
is mixed as
\begin{equation}
f_i = \eta_i\big(0.35\,p^{\mathrm{gb}}_i + 0.60\,\pi_i + 0.05\,\tilde{s}^{\mathrm{ce}}_i\big),
\label{eq:learnedrank}
\end{equation}
where $\rho_i$ is a recency rank, $\Delta_i$ is year distance, and the subject gate $\eta_i$ is $1$ or $0.25$.
The prior, not the trees, carries most of the weight.

The reader is \texttt{deepset/tinyroberta-squad2}.
The span score is the product of the start and end softmax probabilities.
The no-answer score is $1-\sigma(m)$ with the margin $m$ clipped to $[-20,20]$.
A negative margin does not by itself reject the span.
On a multi-sentence passage the reader first sees the date-matched sentence, and a span scoring at least $0.02$ is kept; otherwise the top document is read if its span scores at least $0.05$.
A logistic regression on reader score, no-answer score, and cross-encoder score was fit on 16 labelled rows after standardising the features.
A hand override also keeps a span when the reader score is at least $0.08$ and the cross-encoder score is at least $1$.
This path retrieves 40 candidates and keeps 20, against 80 and 40 for the heuristic systems, so its latency is not comparable even before the model difference.

\section{Experiments}
\label{sec:exp}

\subsection{Systems, and what they do not share}

The frozen tables share the extractive reader, answer projection, sentence selector, and entity or surname cut.
B1--B3 and Adaptive differ on the retrieval decision, the time window, and the sufficiency loop.
\begin{itemize}[leftmargin=1.4em]
  \item \textbf{B1.} The extractive reader with no documents. Hardcoded trivia replies still fire. This is not a closed-book language model.
  \item \textbf{B2.} Always retrieve. Rank by \eqref{eq:blend}. No date window and no sufficiency loop.
  \item \textbf{B3.} Always retrieve. If a year is found, keep candidates whose stored year is within $\pm 2$, and if that set is empty keep the unfiltered pool. Then call the shared ranker without the time-first flag. No sufficiency loop.
  \item \textbf{B4.} Dense candidates union BM25 candidates, scored by an equal blend and then multiplied by $\exp(-|\Delta y|/2)$. No hard cutoff and no sufficiency loop. This is a TempRALM-style decay, not the published TempRALM system.
  \item \textbf{Adaptive.} Router, whitelist skip, adaptive window, and sufficiency loop, under the same reader.
  \item \textbf{Adaptive-Fast.} The saved sufficiency-off ablation. Same answers as Adaptive on the templated and held-out files. Labeled as an ablation-derived variant so the negative loop result stays visible. Its latency is from that later run.
\end{itemize}
The learned path of Section~\ref{sec:learned} stays in the code, flagged experimental.
It is not a row in the frozen comparison.

\subsection{Data}

The templated corpus is generated from timelines of office holders, chief executives, sports winners, awards, and a few laws.
Each fact is a short dated passage whose entity field is the timeline name.
The generator writes questions about the holder at a date, the current holder, and timeless trivia, then fills a target size with paraphrases that prepend a stock phrase to an existing question.
The main file has 803 passages and 3{,}000 questions: 2{,}849 train, 74 validation, and 77 test.
Gold document ids are up to three passages of that entity, nearest the asked date.
One test question asks who was president of Brazil in 2016.
The corpus has no passage for the holder in that year.
Every ranker that stays inside 2016 misses it.
That is a hole in the generated data.

The held-out file rebuilds the same generator with a target of 1{,}200 questions, then moves every question that names one of six entities into the test split.
The six names are every fifth entity in a sorted list: Academy Award for Best Picture, CEO of Apple, CEO of Twitter/X, GDPR, President of Brazil, and UEFA Champions League winners.
Test size becomes 193 because those questions are dumped into test, not because a 15\% split was drawn.
Train and validation of this file do not contain the six names.
The analyzer used at test time is the one trained on the main file, and the main file does contain them.
The split therefore does not hold the entities out of the router.
Almost every other component is a rule and was not trained on question labels.
Withholding names from a training split, while evaluating rules, does not test generalisation of a learned retriever.

The TimeQA subset is built from \texttt{hugosousa/TimeQA}.
The loader would take a half-easy, half-hard sample if rows carried a level field.
We did not record whether that field was present, so this paper does not claim an easy or a hard evaluation.
The built subset has 5{,}931 chunks and 900 questions, 400 train and 500 test.
Chunks are pieces of the page attached to the sampled question, not a retrieval corpus over Wikipedia.
A retrieved chunk is a much easier target than a passage from the open encyclopedia.
Gold ids are chunks containing the answer string, at most eight.

\subsection{Frozen configuration}

All numbers in Section~\ref{sec:results} come from one Checkpoint~B configuration, \texttt{configs/final.yaml}: extractive reader, $k{=}5$, 80 candidates, 40 kept after the cross-encoder, $\alpha=0.6$, confidence and sufficiency thresholds $0.55$, at most two sufficiency repairs, and a 60-second query timeout.
Sentence selection, date and sentence rules, answer projection, and the entity or surname cut are on.
The learned path is off.
\texttt{evaluate.py} loads this file for every system.

Earlier drafts mixed a Checkpoint~A full test with a 16-question Checkpoint~B slice.
Those files are archived under \texttt{artifacts/\_archive\_pre\_fix/} and are not the tables below.

Runs are on a laptop CPU.
No GPU was used.

The reference date inside \eqref{eq:cur} and inside temporal accuracy is 1 October 2025, the date the templated exam was built against.
Current-question behaviour is measured on that templated exam only.

The main analyzer was trained on the templated training questions with six augmentations, 17{,}094 rows, batch size 8.
The saved log contains 7 epochs.
The configured budget was 15.
We do not have a record of why training stopped, and we did not check whether an augmentation of a validation question occurs in the training rows.
Validation necessity accuracy and temporal accuracy are both $1.0$.
On this data the label is often the word ``current'' or a written year.
The figure says the classifier fits those cues.
It does not say the classifier reads time in Wikipedia prose.
No real-text test of the router is reported.

\section{Results}
\label{sec:results}

Tables~\ref{tab:syn77}, \ref{tab:held193}, \ref{tab:timeqa500}, and \ref{tab:ablations} come from the frozen eval directory.
Every system uses the same extractive reader, the same projection, and the same entity or surname cut.

% Auto-generated by scripts/compile_final_results.py. Do not edit by hand.

\begin{table*}[t]
\centering
\caption{Frozen Checkpoint B. Templated test, $n{=}77$. EM is exact match, F1 is token F1, Cont is containment, P@5 and R@5 use gold document ids, Temp is the timestamp metric, Ret is retrieval rate, mean $s$ and med $s$ are seconds.}
\label{tab:syn77}
\small
\begin{tabular}{lccccccccc}
\toprule
System & EM & F1 & Cont & P@5 & R@5 & Temp & Ret & mean $s$ & med $s$ \\
\midrule
B1 no retrieval & 0.130 & 0.167 & 0.130 & 0.000 & 0.000 & 0.247 & 0.000 & 0.000 & 0.000 \\
B2 standard & 0.649 & 0.676 & 0.727 & 0.411 & 0.655 & 0.662 & 1.000 & 1.551 & 1.432 \\
B3 fixed $\pm$2y & 0.935 & 0.951 & 0.987 & 0.608 & 0.943 & 1.000 & 1.000 & 1.555 & 1.452 \\
B4 decay & 0.662 & 0.684 & 0.727 & 0.435 & 0.695 & 0.675 & 1.000 & 0.031 & 0.030 \\
Adaptive & 0.948 & 0.960 & 0.987 & 0.668 & 0.954 & 1.000 & 0.870 & 1.690 & 1.573 \\
Adaptive-Fast & 0.948 & 0.960 & 0.987 & 0.668 & 0.954 & 1.000 & 0.870 & 1.179 & 1.290 \\
\bottomrule
\end{tabular}
\end{table*}

\begin{table*}[t]
\centering
\caption{Frozen Checkpoint B. Questions naming six entities, $n{=}193$. Columns match Table~\ref{tab:syn77}.}
\label{tab:held193}
\small
\begin{tabular}{lccccccccc}
\toprule
System & EM & F1 & Cont & P@5 & R@5 & Temp & Ret & mean $s$ & med $s$ \\
\midrule
B1 no retrieval & 0.000 & 0.002 & 0.000 & 0.000 & 0.000 & 0.000 & 0.000 & 0.000 & 0.000 \\
B2 standard & 0.332 & 0.361 & 0.409 & 0.346 & 0.577 & 0.596 & 1.000 & 1.569 & 1.489 \\
B3 fixed $\pm$2y & 0.684 & 0.696 & 0.710 & 0.591 & 0.879 & 1.000 & 1.000 & 1.546 & 1.410 \\
B4 decay & 0.358 & 0.369 & 0.383 & 0.383 & 0.639 & 0.497 & 1.000 & 0.040 & 0.038 \\
Adaptive & 0.762 & 0.766 & 0.762 & 0.682 & 0.912 & 1.000 & 1.000 & 2.584 & 1.715 \\
Adaptive-Fast & 0.762 & 0.766 & 0.762 & 0.682 & 0.912 & 1.000 & 1.000 & 0.663 & 0.649 \\
\bottomrule
\end{tabular}
\end{table*}

\begin{table*}[t]
\centering
\caption{Frozen Checkpoint B. TimeQA test, $n{=}500$. EM and F1 use a span extracted from the selected sentence; Cont uses the sentence. Temp is omitted.}
\label{tab:timeqa500}
\small
\begin{tabular}{lccccccccc}
\toprule
System & EM & F1 & Cont & P@5 & R@5 & Temp & Ret & mean $s$ & med $s$ \\
\midrule
B1 no retrieval & 0.000 & 0.013 & 0.014 & 0.000 & 0.000 & --- & 0.000 & 0.000 & 0.000 \\
B2 standard & 0.006 & 0.077 & 0.260 & 0.235 & 0.468 & --- & 1.000 & 6.056 & 6.118 \\
B3 fixed $\pm$2y & 0.010 & 0.089 & 0.270 & 0.237 & 0.392 & --- & 1.000 & 6.074 & 6.116 \\
B4 decay & 0.010 & 0.084 & 0.234 & 0.174 & 0.358 & --- & 1.000 & 0.062 & 0.059 \\
Adaptive & 0.008 & 0.088 & 0.296 & 0.241 & 0.488 & --- & 1.000 & 8.736 & 6.492 \\
Adaptive-Fast & 0.008 & 0.088 & 0.298 & 0.243 & 0.491 & --- & 1.000 & 3.184 & 3.160 \\
\bottomrule
\end{tabular}
\end{table*}

\begin{table*}[t]
\centering
\caption{Ablations of Adaptive on the frozen configuration. Each row turns off one component.}
\label{tab:ablations}
\small
\begin{tabular}{llccccc}
\toprule
Data & Ablation & EM & F1 & Cont & Ret & med $s$ \\
\midrule
\multirow{6}{*}{synthetic} & Adaptive (full) & 0.948 & 0.960 & 0.987 & 0.870 & 1.573 \\
 & router off (always retrieve) & 0.948 & 0.960 & 0.987 & 1.000 & 1.456 \\
 & adaptive window off (fixed $\pm$2y) & 0.948 & 0.960 & 0.987 & 0.870 & 1.349 \\
 & sufficiency loop off & 0.948 & 0.960 & 0.987 & 0.870 & 1.290 \\
 & surname/entity cut off & 0.922 & 0.934 & 0.961 & 0.870 & 1.305 \\
 & sentence selector off & 0.948 & 0.960 & 0.987 & 0.870 & 1.276 \\
\midrule
\multirow{6}{*}{heldout} & Adaptive (full) & 0.762 & 0.766 & 0.762 & 1.000 & 1.715 \\
 & router off (always retrieve) & 0.762 & 0.766 & 0.762 & 1.000 & 0.817 \\
 & adaptive window off (fixed $\pm$2y) & 0.762 & 0.766 & 0.762 & 1.000 & 0.689 \\
 & sufficiency loop off & 0.762 & 0.766 & 0.762 & 1.000 & 0.649 \\
 & surname/entity cut off & 0.699 & 0.703 & 0.699 & 1.000 & 0.670 \\
 & sentence selector off & 0.762 & 0.766 & 0.762 & 1.000 & 0.655 \\
\midrule
\multirow{6}{*}{timeqa} & Adaptive (full) & 0.008 & 0.088 & 0.296 & 1.000 & 6.492 \\
 & router off (always retrieve) & 0.008 & 0.088 & 0.296 & 1.000 & 3.185 \\
 & adaptive window off (fixed $\pm$2y) & 0.010 & 0.089 & 0.272 & 1.000 & 3.141 \\
 & sufficiency loop off & 0.008 & 0.088 & 0.298 & 1.000 & 3.160 \\
 & surname/entity cut off & 0.006 & 0.076 & 0.252 & 1.000 & 5.430 \\
 & sentence selector off & 0.002 & 0.048 & 0.168 & 1.000 & 4.300 \\
\bottomrule
\end{tabular}
\end{table*}


\subsection{Templated test}

Table~\ref{tab:syn77} is the full 77-question templated test.
Adaptive containment is $0.987$, the same as B3: 76 of 77 questions.
Exact match is $0.948$ against $0.935$, and token F1 is $0.960$ against $0.951$.
That is one extra exact match after the reader is shared.
Precision at 5 moves from $0.608$ to $0.668$; recall at 5 is $0.954$ against $0.943$.
The remaining retrieval difference is the adaptive window and the sufficiency loop, not an exclusive entity field.

Retrieval rate $0.870$ is 67 of 77 questions.
The other 10 are the whitelist.
Mean latency is $1.690$\,s for Adaptive and $1.555$\,s for B3; medians are $1.573$\,s and $1.452$\,s.
Skipping those 10 questions does not make Adaptive faster.
B2, which has no date logic, is worse on every answer metric (containment $0.727$, exact match $0.649$).
B4, dense plus BM25 with exponential decay and no hard window, lands next to B2: exact match $0.662$, containment $0.727$.
That gap is the comparison a date filter is supposed to win.
It does not separate the adaptive window from the fixed two-year filter.
Adaptive-Fast, the saved sufficiency-off ablation, matches Adaptive on every answer metric here (exact match $0.948$, containment $0.987$).
Its mean latency in that later run is $1.179$\,s, against $1.690$\,s for the four-system Adaptive row.

Temporal accuracy is $1.000$ for both Adaptive and B3.
For B3 that is the circular reading in Section~\ref{sec:setup}.
Paired McNemar tests and 10{,}000-resample bootstrap intervals are in Section~\ref{sec:sig}.
The exact-match difference versus B3 is one question, and the test does not call it a difference ($p=1.000$).

\subsection{Six entity names}

Table~\ref{tab:held193} is the 193-question file.
Adaptive exact match is $0.762$, against $0.684$ for B3 and $0.332$ for B2.
Containment matches exact match for Adaptive at $0.762$, and is $0.710$ and $0.409$ for B3 and B2.
Token F1 follows: $0.766$, $0.696$, $0.361$.
Precision at 5 is $0.682$ against $0.591$ and $0.346$.
Recall at 5 is $0.912$ against $0.879$ and $0.577$.
Here Adaptive does help relative to a fair fixed window.
The paired exact-match difference versus B3 is $+0.078$ with 95\% interval $[+0.041, +0.119]$ and McNemar $p=6.10\times 10^{-5}$ (15 questions only Adaptive, none only B3).
Containment is $+0.052$ $[+0.021, +0.083]$, $p=0.002$.
B4 is not close: exact match $0.358$, containment $0.383$.
Versus B4, exact match differs by $+0.404$ $[+0.316, +0.492]$, $p=3.06\times 10^{-15}$.
Adaptive-Fast matches the full answer scores (exact match $0.762$) with ablation-run mean latency $0.663$\,s against $2.584$\,s for the four-system Adaptive row.
Retrieval rate is $1$.
The ``when'' decision does nothing on this split.
Mean latency is $2.584$\,s against $1.546$\,s for B3; medians are $1.715$\,s and $1.410$\,s.
The mean--median gap is the sufficiency loop on a minority of questions, not a cost paid on every question.

B3's temporal accuracy is again $1.000$, and so is Adaptive's.
Neither figure is answer quality.

\subsection{TimeQA, 500 questions}

Table~\ref{tab:timeqa500} is the full TimeQA test under the frozen config.
Exact match and token F1 use a span extracted from the selected sentence.
Containment still uses the sentence, because a 400-character sentence that contains the gold string is not an exact match.

Adaptive containment is $0.296$, against $0.270$ for B3 and $0.260$ for B2: 148, 135, and 130 of 500 questions.
Exact match is $0.008$ against $0.010$ and $0.006$: four, five, and three questions.
Token F1 is $0.088$ against $0.089$ and $0.077$.
The extractive span does not turn this reader into a TimeQA-paper number.
What Adaptive changes on this file is recall at 5 ($0.488$ against $0.392$ for B3 and $0.468$ for B2) and a small containment gain.
Precision at 5 is $0.241$, $0.237$, and $0.235$.
Retrieval rate is $1$.
Mean latency is $8.736$\,s against about $6.06$\,s; medians are $6.492$\,s against $6.116$\,s.
B3 remains a weak control on this corpus: stored timestamps are often a birth year or 1 January 1970, so a $\pm 2$ year filter can drop the paragraph that names the asked year.
B4 containment is $0.234$ and exact match is $0.010$, the same exact-match count as B3 (five questions) and one ahead of Adaptive.
The containment difference versus B4 is $+0.062$ $[+0.028, +0.096]$, $p=5.39\times 10^{-4}$.
Versus B3 the containment difference is $+0.026$ $[ -0.004, +0.056 ]$, $p=0.117$, so that gain is not a stable finding.
Exact match versus B3 is $-0.002$, $p=1.000$.
Adaptive-Fast containment is $0.298$ against $0.296$, exact match unchanged at $0.008$, and the ablation-run mean is $3.184$\,s against $8.736$\,s.
The learned path is not in this table.

\subsection{Paired tests}
\label{sec:sig}

Every interval below is a paired bootstrap of Adaptive minus the named system, 10{,}000 resamples, seed 0.
Every $p$-value is an exact two-sided McNemar test on discordant questions.
The full grid, including B2, is \texttt{artifacts/significance.json}.

On the templated test, Adaptive versus B3 containment is a tie ($p=1.000$) and exact match is one question ($p=1.000$).
Versus B4, exact match is $+0.286$ $[+0.182, +0.390]$, $p=2.98\times 10^{-6}$.
On the held-out file the entity-cut ablation is the answer change that survives a paired test: exact match $+0.062$ $[+0.031, +0.098]$, $p=4.88\times 10^{-4}$, 12 questions only the full system.
On TimeQA, sentence-selector off drops containment by $0.128$ $[+0.084, +0.172]$, $p=2.89\times 10^{-8}$.
Entity-cut off drops containment by $0.044$ $[+0.024, +0.066]$, $p=1.13\times 10^{-4}$.
Window off is $+0.024$ $[ -0.004, +0.052 ]$, $p=0.134$.
Sufficiency off is a tie on exact match and one extra contained answer for the ablated system ($0.298$ against $0.296$, $p=1.000$).
The loop is a negative result: it does not buy answers, and the full Adaptive row is the slow one.

\subsection{Ablations}

Table~\ref{tab:ablations} turns off one Adaptive component at a time on each full test: the router (always retrieve), the adaptive window (fixed $\pm 2$ years), the sufficiency loop, the surname or entity cut, and the sentence selector.
The other frozen settings stay on.
Answer scores in this table are comparable.
Latency is from a later Adaptive-only run, so it is not matched to the four-system means in Tables~\ref{tab:syn77}--\ref{tab:timeqa500}.

On the templated and held-out files, turning off the router, the adaptive window, the sufficiency loop, or the sentence selector leaves exact match and containment unchanged.
The entity cut is the component that moves those answers: synthetic exact match falls from $0.948$ to $0.922$, and held-out exact match from $0.762$ to $0.699$.
Sufficiency off is faster and does not change those scores.

On TimeQA the sentence selector is the largest drop: containment $0.168$ against $0.296$, exact match $0.002$ against $0.008$.
The entity cut also matters (containment $0.252$).
Turning the adaptive window into a fixed $\pm 2$ year filter lowers containment to $0.272$ and recall at 5 to $0.397$, which is the B3 pattern.
Sufficiency off is $0.298$ containment against $0.296$: the loop does not buy answers on this file and it costs time.
The router ablation matches the full system, because retrieval rate is already $1$.

\section{What these measurements do not show}
\label{sec:limits}

The sentences below are the limits of this draft.
They are not future work written as a compliment to the method.

The frozen comparison shares the reader, the projection, and the entity cut.
What remains different is skip versus always-retrieve, the window, and the sufficiency loop.
Ablations in Table~\ref{tab:ablations} turn those pieces off one at a time.
B4 is a TempRALM-style decay ranker written for this comparison.
It is not the authors' TempRALM system, and Adaptive-RAG is still not reimplemented.

B3 is a reasonable control when the stored timestamp is the fact date.
It is a weak control on TimeQA, and its perfect temporal accuracy is largely the filter scoring itself.

The skip decision was not tested against a model that knows the trivia without the whitelist.
Latency does not improve.
Held-out retrieval rate is $1$.
The project title's first half is not a result on that split.

Constants throughout Section~\ref{sec:method} --- the blend weights, $\lambda$, $\alpha$, the window table, the 45-year cutoff, the 15 June midpoint, the $0.55$ threshold, the $0.01$ tie --- were chosen by hand.
Some of the TimeQA rules were chosen while looking at development failures.
There is no untouched development split for those choices.

The learned path is a frozen reader, a tree model on 357 rows, and a logistic model on 16 rows plus an override.
It is experimental and is not in the frozen tables.
Its router accuracy was computed on the split used to pick $T$.

Containment is bidirectional and, for this reader, close to ``did the gold string occur in the chosen sentence?''
Temporal accuracy does not read the answer, and the same words are used for the router's class accuracy in Section~\ref{sec:learned}.
The two quantities are not the same.

The current-question branch is scored against a frozen day on templated text.
FreshQA and RealTime QA were not run.
A deployment clock and this exam clock are different.

We do not release a public code URL with this draft.
The implementation is the project tree that produced the tables.

\section{Conclusion}

The pipeline can prefer a dated sentence and can keep a subject-matched window.
On the templated test, with a shared reader, that is worth one extra exact match over a fixed two-year filter and nothing on containment.
Precision at 5 is higher.
Skipping retrieval is a trivia whitelist, used on 10 of 77 questions, unused on the 193-question split, and invisible in the latency.
On the six-entity split the same frozen system is ahead of the fair fixed window on exact match, containment, and precision at 5, and it is slower.
On 500 TimeQA questions the containment edge over B3 is not significant ($p=0.117$), the edge over B4 is ($p=5.39\times 10^{-4}$), and exact match is not an Adaptive win.
Turning the sufficiency loop off does not hurt those scores and it removes the slow tail.
That negative result is part of the finding, not a footnote to hide.

\begin{thebibliography}{99}

\bibitem{lewis2020rag}
P.~Lewis, E.~Perez, A.~Piktus, F.~Petroni, V.~Karpukhin, N.~Goyal, H.~K\"uttler, M.~Lewis, W.~Yih, T.~Rockt\"aschel, S.~Riedel, and D.~Kiela.
Retrieval-augmented generation for knowledge-intensive NLP tasks.
In \emph{Advances in Neural Information Processing Systems}, 2020.

\bibitem{karpukhin2020dpr}
V.~Karpukhin, B.~Oguz, S.~Min, P.~Lewis, L.~Wu, S.~Edunov, D.~Chen, and W.~Yih.
Dense passage retrieval for open-domain question answering.
In \emph{Proceedings of EMNLP}, 2020.

\bibitem{izacard2023atlas}
G.~Izacard, P.~Lewis, M.~Lomeli, L.~Hosseini, F.~Petroni, T.~Schick, J.~Dwivedi-Yu, A.~Joulin, S.~Riedel, and E.~Grave.
Atlas: Few-shot learning with retrieval augmented language models.
\emph{Journal of Machine Learning Research}, 24(251):1--43, 2023.

\bibitem{asai2024selfrag}
A.~Asai, Z.~Wu, Y.~Wang, A.~Sil, and H.~Hajishirzi.
Self-RAG: Learning to retrieve, generate, and critique through self-reflection.
In \emph{International Conference on Learning Representations}, 2024.

\bibitem{jeong2024adaptiverag}
S.~Jeong, J.~Baek, S.~Cho, S.~J.~Hwang, and J.~C.~Park.
Adaptive-RAG: Learning to adapt retrieval-augmented large language models through question complexity.
In \emph{Proceedings of NAACL}, 2024.

\bibitem{li2003time}
X.~Li and W.~B.~Croft.
Time-based language models.
In \emph{Proceedings of CIKM}, pages 469--475, 2003.

\bibitem{kanhabua2011ranking}
N.~Kanhabua and K.~N{\o}rv{\aa}g.
A comparison of time-aware ranking methods.
In \emph{Proceedings of SIGIR}, 2011.

\bibitem{jia2018tempquestions}
Z.~Jia, A.~Abujabal, R.~Saha~Roy, J.~Str\"otgen, and G.~Weikum.
TempQuestions: A benchmark for temporal question answering.
In \emph{Companion of The Web Conference}, 2018.

\bibitem{chen2021timeqa}
W.~Chen, X.~Wang, and W.~Y.~Wang.
A dataset for answering time-sensitive questions.
In \emph{NeurIPS Datasets and Benchmarks}, 2021.
arXiv:2108.06314.

\bibitem{zhang2021situatedqa}
M.~J.~Q.~Zhang and E.~Choi.
SituatedQA: Incorporating extra-linguistic contexts into QA.
In \emph{Proceedings of EMNLP}, 2021.

\bibitem{liska2022streamingqa}
A.~Li\v{s}ka, T.~Ko\v{c}isk\'y, E.~Gribovskaya, T.~Terzi, E.~Sezener, D.~Agrawal, C.~de~Masson~d'Autume, T.~Scholtes, M.~Zaheer, S.~Young, E.~Gilsenan-McMahon, S.~Austin, P.~Blunsom, and A.~Lazaridou.
StreamingQA: A benchmark for adaptation to new knowledge over time in question answering models.
In \emph{Proceedings of ICML}, 2022.

\bibitem{kasai2023realtimeqa}
J.~Kasai, K.~Sakaguchi, Y.~Takahashi, R.~Le~Bras, A.~Asai, X.~Yu, D.~Radev, N.~A.~Smith, Y.~Choi, and K.~Inui.
RealTime QA: What's the answer right now?
In \emph{NeurIPS Datasets and Benchmarks}, 2023.

\bibitem{dhingra2022talm}
B.~Dhingra, J.~R.~Cole, J.~M.~Eisenschlos, D.~Gillick, J.~Eisenstein, and W.~W.~Cohen.
Time-aware language models as temporal knowledge bases.
\emph{Transactions of the Association for Computational Linguistics}, 10:257--273, 2022.

\bibitem{vu2023freshqa}
T.~Vu, M.~Iyyer, X.~Wang, N.~Constant, Jerry~Wei, Jason~Wei, C.~Tar, Y.-H.~Sung, D.~Zhou, Q.~Le, and T.~Luong.
FreshLLMs: Refreshing large language models with search engine augmentation.
arXiv:2310.03214, 2023.

\bibitem{gade2024tempralm}
A.~Gade and J.~Jetcheva.
It's about time: Incorporating temporality in retrieval augmented language models.
arXiv:2401.13222, 2024.

\bibitem{rajpurkar2018squad2}
P.~Rajpurkar, R.~Jia, and P.~Liang.
Know what you don't know: Unanswerable questions for SQuAD.
In \emph{Proceedings of ACL}, 2018.

\bibitem{xiao2024cpack}
S.~Xiao, Z.~Liu, P.~Zhang, N.~Muennighoff, D.~Lian, and J.-Y.~Nie.
C-Pack: Packaged resources to advance general Chinese embedding.
In \emph{Proceedings of SIGIR}, 2024.
arXiv:2309.07597.

\bibitem{johnson2019faiss}
J.~Johnson, M.~Douze, and H.~J\'egou.
Billion-scale similarity search with GPUs.
\emph{IEEE Transactions on Big Data}, 7(3):535--547, 2019.

\bibitem{wang2020minilm}
W.~Wang, F.~Wei, L.~Dong, H.~Bao, N.~Yang, and M.~Zhou.
MiniLM: Deep self-attention distillation for task-agnostic compression of pre-trained transformers.
In \emph{International Conference on Learning Representations}, 2020.

\bibitem{nguyen2016msmarco}
T.~Nguyen, M.~Rosenberg, X.~Song, J.~Gao, S.~Tiwary, R.~Majumder, and L.~Deng.
MS MARCO: A human generated machine reading comprehension dataset.
In \emph{CoCo@NeurIPS}, 2016.

\bibitem{reimers2019sbert}
N.~Reimers and I.~Gurevych.
Sentence-BERT: Sentence embeddings using Siamese BERT-networks.
In \emph{Proceedings of EMNLP}, 2019.

\bibitem{sanh2019distilbert}
V.~Sanh, L.~Debut, J.~Chaumond, and T.~Wolf.
DistilBERT, a distilled version of BERT: Smaller, faster, cheaper and lighter.
arXiv:1910.01108, 2019.

\bibitem{devlin2019bert}
J.~Devlin, M.-W.~Chang, K.~Lee, and K.~Toutanova.
BERT: Pre-training of deep bidirectional transformers for language understanding.
In \emph{Proceedings of NAACL}, 2019.

\bibitem{friedman2001gbm}
J.~H.~Friedman.
Greedy function approximation: A gradient boosting machine.
\emph{Annals of Statistics}, 29(5):1189--1232, 2001.

\end{thebibliography}

\end{document}
